% Chapter 03 — Development Workflow & Tools
\section{Chương 3: Phát triển và công cụ}
\ChapterDivider{Chương 3}{Phát triển và công cụ}{Từ mã nguồn đến quan sát được sự kiện thật}

\begin{frame}{Quy trình phát triển eBPF (ví dụ bằng C)}
  \begin{columns}[T,onlytextwidth]
    \begin{column}{0.49\textwidth}
      \cue{1}{\textbf{Write}: viết chương trình eBPF bằng C.}
      \vspace{0.18cm}
      \cue{2}{\textbf{Compile}: Clang/LLVM tạo eBPF bytecode trong tệp object.}
      \vspace{0.18cm}
      \cue{3}{\textbf{Load + Verify}: loader yêu cầu kernel nạp; Verifier kiểm tra.}
    \end{column}
    \begin{column}{0.49\textwidth}
      \cue{4}{\textbf{Attach}: gắn program vào hook tương thích.}
      \vspace{0.18cm}
      \cue{5}{\textbf{Execute}: sự kiện phù hợp kích hoạt program.}
      \vspace{0.18cm}
      \cue{6}{\textbf{Collect}: đọc map hoặc nhận các sự kiện ở User Space.}
    \end{column}
  \end{columns}
  \vspace{0.15cm}
  \begin{block}{Hai điểm dễ nhầm}
    \textbf{Load} không đồng nghĩa với \textbf{Attach}. JIT là tùy chọn;
    công cụ cấp cao có thể tự xử lý nhiều bước thay người dùng.
  \end{block}
\end{frame}

\begin{frame}{Các công cụ eBPF dùng để làm gì?}
  \begin{columns}[T,onlytextwidth]
    \begin{column}{0.49\textwidth}
      \begin{block}{\texttt{bpftrace} — tracing nhanh}
        Viết và chạy các phép theo dõi bằng cú pháp ngắn; phù hợp để
        minh họa sự kiện trong demo.
      \end{block}
      \begin{block}{\texttt{bpftool} — kiểm tra hệ thống}
        Xem tính năng BPF được kernel hỗ trợ, quản lý program, map, link.
      \end{block}
    \end{column}
    \begin{column}{0.49\textwidth}
      \begin{block}{\texttt{libbpf} — phát triển ứng dụng C}
        Thư viện hỗ trợ nạp, gắn và quản lý chương trình eBPF.
      \end{block}
      \begin{block}{BCC — bộ công cụ tracing}
        Công cụ và thư viện để quan sát, phân tích hành vi và hiệu năng hệ thống.
      \end{block}
    \end{column}
  \end{columns}
  \vspace{0.12cm}
  \begin{alertblock}{Lựa chọn cho phần thực hành}
    \texttt{bpftool} để kiểm tra môi trường; \texttt{bpftrace} để quan sát
    system call và đếm sự kiện mà chưa cần tự xây dựng loader bằng C.
  \end{alertblock}
\end{frame}

\begin{frame}{Ví dụ xuyên suốt: theo dõi lời gọi \texttt{openat()}}
  \begin{block}{Ta muốn quan sát điều gì?}
    Đếm số lần các chương trình gọi \texttt{openat()} để yêu cầu mở tệp
    hoặc thư mục, rồi xem thống kê theo tên chương trình.
  \end{block}
  \vspace{0.14cm}
  \begin{center}
    \small
    \fbox{\parbox{0.16\textwidth}{\centering Application\\\texttt{openat()}}}
    $\rightarrow$
    \fbox{\parbox{0.24\textwidth}{\centering Tracepoint\\\texttt{sys\_enter\_openat}}}
    $\rightarrow$
    \fbox{\parbox{0.17\textwidth}{\centering eBPF\\Program}}
    $\rightarrow$
    \fbox{\parbox{0.12\textwidth}{\centering Map\\count}}
  \end{center}
  \begin{center}
    $\longrightarrow$ \textbf{bpftrace đọc map và hiển thị thống kê theo \texttt{comm}}
  \end{center}
  \vspace{0.1cm}
  \cue{1}{\texttt{comm} là tên tác vụ/tiến trình do kernel cung cấp, \textbf{không phải} PID.}
  \cue{2}{Tracepoint ở \textit{entry} ghi nhận lời gọi, \textbf{không} xác nhận mở tệp thành công.}
\end{frame}
